% !TEX root = ../main.tex
% Figure: Mode splitting due to perturbation in coupled resonators
% Chapter: ch11 - CMT and Perturbation Theory for Mode Splitting
% Description: Degenerate modes split when symmetry is broken

\begin{tikzpicture}[scale=1.2]
    % Left: Unperturbed degenerate system
    \begin{scope}[shift={(-5,0)}]
      % Two identical rings (resonators)
      \draw[thick, blue] (-1,0) circle (1);
      \draw[thick, red] (1,0) circle (1);
      
      % Overlap region (coupling)
      \fill[green!20] (0,-0.3) ellipse (0.3 and 0.6);
      
      % Field distribution in ring 1
      \draw[very thick, blue, domain=-180:-90, samples=30] 
        plot({-1+cos(\x)+0.2*sin(3*\x)},{sin(\x)});
      
      % Field distribution in ring 2
      \draw[very thick, red, domain=0:90, samples=30] 
        plot({1+cos(\x)-0.2*sin(3*\x)},{sin(\x)});
      
      % Energy level diagram above
      \draw[thick] (-1.5,2) -- (-0.5,2) node[midway, above] {$\omega_0$};
      \draw[thick] (0.5,2) -- (1.5,2) node[midway, above] {$\omega_0$};
      
      \node[above, font=\bfseries] at (0,2.8) {(a) 未微扰：简并态};
      \node[below, align=center] at (0,-1.5) {两个相同谐振腔\\弱耦合，频率简并};
      
      % Coupling strength indication
      \draw[<->, thick, green!60!black] (-0.3,0) -- (0.3,0);
      \node[above, green!60!black] at (0,0.2) {\scriptsize $\kappa$};
    \end{scope}
    
    % Right: Perturbed system with splitting
    \begin{scope}[shift={(2,0)}]
      % Asymmetric rings (one slightly deformed)
      \draw[thick, blue] (-1,0) circle (1);
      \draw[thick, red, rotate around={15:(1,0)}] (1,0) circle (0.95);
      
      % Reduced overlap
      \fill[orange!20] (0.2,-0.2) ellipse (0.2 and 0.4);
      
      % Modified field patterns
      \draw[very thick, blue, domain=-180:-90, samples=30] 
        plot({-1+cos(\x)+0.15*sin(2*\x)},{sin(\x)});
      
      \draw[very thick, red, domain=-30:105, samples=30] 
        plot({1+0.95*cos(\x)-0.25*sin(2*\x)},{0.95*sin(\x)});
      
      % Energy level splitting
      \draw[thick] (-1.5,2.3) -- (-0.5,2.3) node[midway, above] {$\omega_+$};
      \draw[thick] (-1.5,1.7) -- (-0.5,1.7) node[midway, below] {$\omega_-$};
      
      \draw[thick] (0.5,2.5) -- (1.5,2.5) node[midway, above] {$\omega_0+\Delta$};
      \draw[thick] (0.5,1.5) -- (1.5,1.5) node[midway, below] {$\omega_0-\Delta$};
      
      \node[above, font=\bfseries] at (0,2.8) {(b) 微扰后：简并解除};
      \node[below, align=center] at (0,-1.5) {对称性破缺\\能级分裂$2\Delta$};
      
      % Splitting magnitude
      \draw[<->, very thick, purple] (1.7,1.5) -- (1.7,2.5);
      \node[right, purple, font=\bfseries] at (1.7,2) {$2\Delta$};
    \end{scope}
    
    % Center: Transition arrow
    \draw[->, ultra thick, orange, bend left=20] (-2,1.5) to[out=30,in=150] 
      node[above, align=center, font=\small] {施加微扰\\形变/折射率变化/应力} (0,1.5);
    
    % Bottom: Mathematical formulation
    \node[draw, align=center, font=\small] at (0,-2.5) {
      微扰哈密顿量矩阵元:\\
      $H'_{mn}=\langle\psi_m^{(0)}|\delta\varepsilon|\psi_n^{(0)}\rangle$\\[4pt]
      本征值分裂:$\Delta\omega = \frac{\omega_0}{2}\sqrt{\left(\frac{\delta\varepsilon}{\varepsilon}\right)^2 + 4\kappa^2}$
    };
    
    % Physical examples sidebar
    \node[align=left, font=\footnotesize, anchor=south east] at (-3.5,-3) {
      \textbf{实际微扰来源}:\\
      $\bullet$ 制造误差（孔径偏差$\pm 2\,\mathrm{nm}$）\\
      $\bullet$ 热梯度引起的折射率不均匀\\
      $\bullet$ 机械应力导致的双折射\\
      $\bullet$ 载流子注入的空间非均匀性
    };
  \end{tikzpicture}
